\section{Model, solution framework, and sampling laws}
\label{sec:model}

The size variable $x\in E:=(0,\infty)$ is proportional to particle mass.
We fix its reference unit throughout, so that logarithms of size and sums such
as $1+x$ below are dimensionless. The finite nonnegative measure $n_t$ describes
particle number per unit volume. For $p\geq0$, write
\begin{equation}
 M_p(t)=\int_E x^p n_t(\dd x),\qquad N(t)=M_0(t),\qquad m=M_1(0).
 \label{eq:moments}
\end{equation}
We assume $0<N_0:=N(0)<\infty$ and $0<m<\infty$.

\begin{assumption}[Rates and daughters]\label{ass:rates}
The functions $\lambda,\sigma:[0,\infty)\to[0,\infty)$ are measurable and
locally integrable. Coagulation has the additive kernel
$K_t(x,y)=\lambda(t)(x+y)$. Fragmentation has per-particle selection rate
$\sigma(t)$. Its expected daughter measure is a measurable kernel
$(t,x)\mapsto b_{t,x}$ on $E$, satisfying
\begin{equation}
 b_{t,x}(E\setminus(0,x))=0,\qquad
 \int_E b_{t,x}(\dd z)=2,\qquad
 \int_E z\,b_{t,x}(\dd z)=x.
 \label{eq:daughters}
\end{equation}
These conditions hold for every $(t,x)$; the values of the model at a common
null set of times may be chosen arbitrarily subject to these conditions.
\end{assumption}

Only the expected daughter count and mass enter the deterministic equation.
An actual binary split has expected measure
$\E[\delta_{Vx}+\delta_{(1-V)x}]$ for a random $V\in(0,1)$.
Conditions \eqref{eq:daughters} allow a larger class of expected measures:
they do not require symmetry about $x/2$, which an eventwise complementary
binary split would impose. Results about a finite system of particles will
explicitly require eventwise binary mass conservation. Unless stated otherwise,
the deterministic results use only \eqref{eq:daughters}.

Write $\pair{f}{\mu}=\int_E f\dd\mu$ whenever this integral is defined.
For a Borel test $f$, define
\begin{equation}
 \Delta f(x,y)=f(x+y)-f(x)-f(y),\qquad
 \mathcal F_t f(x)=\int_E f(z)b_{t,x}(\dd z)-f(x).
 \label{eq:increments}
\end{equation}
The weak population balance is
\begin{align}
 \pair{f}{n_t}-\pair{f}{n_0}
 &=\frac12\int_0^t\lambda(s)\iint_{E^2}(x+y)\Delta f(x,y)
                  n_s(\dd x)n_s(\dd y)\dd s\notag\\
 &\quad+\int_0^t\sigma(s)\int_E\mathcal F_s f(x)n_s(\dd x)\dd s.
 \label{eq:weak}
\end{align}
The factor $1/2$ counts each unordered coagulation pair once.

\begin{definition}[Mass-conserving solution]\label{def:solution}
A solution is a curve of nonnegative finite measures $n_t$ on $E$ such that
$t\mapsto n_t(A)$ is measurable for each Borel set $A\subset E$,
with $\sup_{s\leq T}N(s)<\infty$ for every $T<\infty$, constant first moment
$M_1(t)=m$, and \eqref{eq:weak} for every bounded Borel $f$ and every $t\geq0$.
\end{definition}

All terms in this definition are absolutely integrable on bounded time
intervals: their absolute values are bounded using
$\iint(x+y)n_s(\dd x)n_s(\dd y)=2mN(s)$ and
$|\mathcal F_s f|\leq3\norm{f}_\infty$.
In particular the solution is continuous in total variation.
For a signed measure $\mu$, write
\begin{equation}
 \norm{\mu}_w=\int_E(1+x)|\mu|(\dd x),\qquad
 \norm{\mu}_{\mathrm{var}}=|\mu|(E).
 \label{eq:variation}
\end{equation}
For probability measures the distance
$\norm{\mu-\nu}_{\TV}:=\sup_A|\mu(A)-\nu(A)|$
equals $\tfrac12\norm{\mu-\nu}_{\mathrm{var}}$.

\begin{theorem}[A well-posed solution class]\label{thm:existence}
Under \cref{ass:rates}, suppose additionally that $M_2(0)<\infty$.
There is a unique mass-conserving solution with locally bounded second
moment. It is continuous in the weighted-variation norm $\norm{\cdot}_w$
and satisfies
\begin{equation}
 M_2(t)\leq M_2(0)\exp\left(2m\int_0^t\lambda(s)\dd s\right).
 \label{eq:M2bound}
\end{equation}
No continuity of the daughter kernel in the parent variable is required.
\end{theorem}

A self-contained proof appears in \cref{app:existence}. It uses a
weighted-variation stability estimate and approximations that converge in
that norm. This avoids passing a merely measurable parent-dependent kernel
through a narrow limit. The proof also avoids treating weak integrals of
moving atomic measures as Bochner integrals. We use the second moment as a
sufficient condition for this construction, not as a necessary condition for
the population balance or the results below. In the time-independent
selfsimilar setting, weaker initial-moment assumptions are available for
fixed relative-fragment laws satisfying the dislocation hypotheses of
\citet{Cepeda2014}. That work treats homogeneous-like coagulation and
multiple fragmentation using a Wasserstein-type distance. The finite first
moment is the relevant homogeneity moment for the additive kernel. Neither
that result nor the proof here is a claim of novelty for existence methods.

The estimates that follow require only \cref{def:solution}, unless a theorem
explicitly adds moment assumptions. Thus they also apply to any separately
constructed finite-count, finite-mass solution outside \cref{thm:existence}.
This distinction is useful for logarithmic observables: neither a finite
second size moment nor finite count and mass implies a finite negative
logarithmic moment.

Set
\begin{equation}
 b(t)=m\lambda(t),\qquad
 \Bc(t)=\int_0^t b(s)\dd s,\qquad
 F(t)=\int_0^t\sigma(s)\dd s,\qquad C(t)=\Bc(t)+F(t).
 \label{eq:clocks}
\end{equation}
Here $b(t)$ is a scalar rate; the two-index symbol $b_{t,x}$ always denotes
the daughter measure. Testing \eqref{eq:weak} with $f=1$ shows that $N$ is
locally absolutely continuous, with
\begin{equation}
 N'(t)=(\sigma(t)-b(t))N(t)\quad\text{for almost every }t,\qquad
 N(t)=N_0e^{F(t)-\Bc(t)}\quad(t\geq0).
 \label{eq:count}
\end{equation}
In particular $N(t)>0$ at each finite time. The number and mass sampling laws
are respectively
\begin{equation}
 \eta_t(\dd x)=\frac{n_t(\dd x)}{N(t)},\qquad
 \pi_t(\dd x)=\frac{x\,n_t(\dd x)}{m},\qquad
 \frac{\dd\pi_t}{\dd\eta_t}(x)=\frac{N(t)x}{m}.
 \label{eq:sampling}
\end{equation}
The first law samples particles uniformly by number; the second samples a
uniform unit of material and records the size of its containing particle.
Their distinction is central even when their underlying population is the
same. Probabilistic representations of mass sampling in pure coagulation
are established; see, for example, \citet{DeaconuFournierTanre2002}.

For $0<p<1$, define the normalized fractional moment
\begin{equation}
 \mathcal A_p(t)=\frac{M_p(t)}{N(t)^{1-p}m^p}
 =\int_E\left(\frac{\dd\pi_t}{\dd\eta_t}\right)^p\dd\eta_t.
 \label{eq:affinity}
\end{equation}
H\"older's inequality gives $0<\mathcal A_p\leq1$.
At $p=1/2$ it is the Hellinger affinity of the two sampling laws.
No density of $n_t$ with respect to Lebesgue measure is assumed anywhere in
these definitions.
